\section{UK research stream B: Economic outcomes after inheritance}
\label{sec:outcomes}

\subsection{Common data requirement}
To measure what recipients do after inheriting, we need the transfer's \textbf{date, value, source and asset type}; the recipient's prior wealth, debt, income and family background; and repeated outcomes before and after receipt. ELSA, WAS and Understanding Society may support distinct within-survey analyses, subject to variable overlap. Detailed business, earnings or credit outcomes may need separately approved administrative linkage.

\subsection{Candidate research questions}
\begin{description}[style=nextline]
\item[B1. Entrepreneurship]
Does inheritance precede business entry, self-employment or expansion, especially for low-liquid-wealth recipients? Need firm start date, pre-transfer employment/wealth, firm performance and survival. Distinguish starting a business from inheriting a family firm. Possible sources: surveys, HMRC business records and Companies House (linkage conditional).
\item[B2. Housing]
Does it accelerate first home purchase, improve deposits, change housing tenure or reduce mortgages? Need transfer date, property transactions/tenure, mortgage balance and baseline wealth. Start with WAS; Land Registry and lender data would be conditional extensions.
\item[B3. Saving, debt and portfolios]
How are financial, housing and pension assets and secured/unsecured debts rebalanced after receipt? Need repeated gross asset and liability categories, not just net worth. A self-reported use of inheritance is not a complete accounting of a fungible transfer.
\item[B4. Work and retirement]
Do recipients reduce hours, retire early, change jobs or invest in training? Need longitudinal earnings, hours, employment, occupations and pensions. ELSA is relevant later in life; Understanding Society and approved tax data may cover younger workers.
\item[B5. Consumption and resilience]
Does receipt reduce arrears/financial hardship or finance major purchases? Need spending or transaction data and household hardship indicators. Survey spending categories alone may be insufficient for comprehensive consumption estimates.
\item[B6. Family and location]
Does inheritance affect household formation, geographic mobility or subsequent gifts to children? Need household histories, residence and transfer dates, with local trends as controls.
\end{description}

For every question, examine heterogeneity by recipient age, transfer amount, prior liquidity, parental resources, education, gender, siblings and region. These may be distinctive economic mechanisms even without a full causal design.

\subsection{Empirical sequence and cautions}
Start by counting recipients with an observed pre-transfer and post-transfer panel. Plot outcomes in event time, adjust for survey weights/attrition and check anticipation. Bereavement, care responsibilities and selection into receipt can independently affect outcomes. Event plots are descriptive unless an appropriate identification strategy is established; a parental death is not automatically an exogenous instrument for inheritance.

The first accessible results are most likely wealth composition, housing and employment/retirement trajectories. Firm survival, detailed consumption, tax earnings and credit require richer data and permissions.

\subsection{To do}
\begin{itemize}[nosep]
\item Build an outcome-by-dataset matrix covering exact variables, waves and sample overlap.
\item Prioritise feasible WAS/ELSA panels for housing, wealth/debt and employment.
\item Define transfer event windows, receipt amounts, baseline characteristics and comparison samples.
\item Produce simple pre/post descriptive plots with attrition and bereavement checks.
\item Scope more ambitious firm, earnings and credit linkage separately; identify which additional fields are indispensable.
\end{itemize}
